Let's define a mechanism for adding one specific vector to a flagged one.
This technique works even when there are more than one flagged vector.

The flag used for the vectors should have $f_{\ON} = 1$ and $f_{\OFF} = 0$.
We will assume the flag is in the last coordinate.

Let's call $v$ the vector that we want to add. 
We will do it with one $\FF$ layer. 
We will construct a layer such that $FF(h) = \mathbb{I}(h \text{ is flagged})*v$.

\begin{align*}
    & W_1 = id
    & W_2  = \left(\begin{matrix}
        & 0         &\dots  &0      & v_1       \\
        & \vdots    &\ddots &\vdots & \vdots    \\
        & 0         &\dots  &0      & v_d       \\
    \end{matrix}\right) &
    & b_1 = b_2 = \left(\begin{matrix}
        & 0      \\
        & \vdots \\
        & 0      \\
    \end{matrix}\right) 
\end{align*}

Notice that if the coordinate for the flag had been the $k$th instead of the last one, then $v$ would have appeared in the $k$th column of $W_2$ instead of the last.

\medskip

An interesting use of this technique is that allow us to transform any vector in whatever we want, i.e. we can apply constant functions.
This is done exploiting the finite representation system.
We can saturate the vector to the $\maxrep$ and then make it $0$ to finally add our objective.


\[
    h \leftarrow \left(\left(\left( h + 
    \left(\begin{matrix}
    \infty \\
    \vdots \\
    \infty \\
\end{matrix}\right)\right)
    + \left(\begin{matrix}
    \infty \\
    \vdots \\
    \infty \\
\end{matrix}\right)\right)
    + \left(\begin{matrix}
    -\infty \\
    \vdots \\
    -\infty \\
\end{matrix}\right)\right)
    + v
\]
